\documentclass[conference]{IEEEtran}
\IEEEoverridecommandlockouts
% The preceding line is only needed to identify funding in the first footnote. If that is unneeded, please comment it out.
\usepackage{cite}
\usepackage{amsmath,amssymb,amsfonts}

% Algorithms / pseudocode (algorithm2e)
\usepackage[ruled,vlined,linesnumbered]{algorithm2e}
\DontPrintSemicolon
\SetAlgoLined
\SetAlFnt{\footnotesize} % reduce line wrapping risk in IEEE two-column layout
\SetKwInOut{Input}{Input}
\SetKwInOut{Output}{Output}
\SetKw{Continue}{continue}

% Force figures to actually render (avoid draft placeholders)
\usepackage[final]{graphicx}
\usepackage{textcomp}
\usepackage{xcolor}
\usepackage[final]{svg}
% Extra safety: ensure graphicx isn't in draft mode anywhere
\setkeys{Gin}{draft=false}

% Float spacing (tighter, while keeping IEEEtran reasonably readable)
\setlength{\textfloatsep}{7pt plus 1pt minus 2pt}
\setlength{\floatsep}{6pt plus 1pt minus 2pt}
\setlength{\intextsep}{6pt plus 1pt minus 2pt}
\setlength{\abovecaptionskip}{2pt}
\setlength{\belowcaptionskip}{0pt}

% Tables
\usepackage{booktabs} % \toprule, \midrule, \bottomrule
\usepackage{tabularx} % automatic line-wrapping columns
\usepackage{array}    % column modifiers (e.g., >{\raggedright\arraybackslash}X)
\usepackage{threeparttable} % table notes with width matching the tabular block
\usepackage{capt-of} % \captionof{table}{...} inside minipage (side-by-side tables within \columnwidth)
\usepackage{dblfloatfix} % improve placement of two-column floats; can enable bottom placement
\usepackage[section]{placeins} % keep floats from wandering across section boundaries

% Verdict-matrix markers (rounded color tags)
\usepackage{tikz}
\newcommand{\pill}[2]{%
  \tikz[baseline=(n.base)]\node (n) [rounded corners=3pt, fill=#1, inner xsep=5pt, inner ysep=2.2pt]{\bfseries\footnotesize #2};%
}
\newcommand{\yes}{\pill{green!18}{yes}}
\newcommand{\no}{\pill{red!18}{no}}
\newcommand{\nc}{\pill{orange!20}{n/c}}
\newcommand{\NA}{\pill{gray!20}{n/a}}

\usepackage{biblatex}
\addbibresource{refs.bib}
\def\BibTeX{{\rm B\kern-.05em{\sc i\kern-.025em b}\kern-.08em
    T\kern-.1667em\lower.7ex\hbox{E}\kern-.125emX}}
\begin{document}

\title{AC-GATE: Entity-Conditioned Lags\\
{\footnotesize \textsuperscript{*}Auditing Heterogeneous Effective Lags with Entity Proxies}
}



\author{\IEEEauthorblockN{1\textsuperscript{st} Andi Xu}
\IEEEauthorblockA{\textit{School of engineering} \\
\textit{jönköping university}\\
jönköping, Sweden \\
xuan25aw@student.ju.se}
\and
\IEEEauthorblockN{2\textsuperscript{nd} Given Name Surname}
\IEEEauthorblockA{\textit{dept. name of organization (of Aff.)} \\
\textit{name of organization (of Aff.)}\\
City, Country \\
email address or ORCID}
\and
\IEEEauthorblockN{3\textsuperscript{rd} Given Name Surname}
\IEEEauthorblockA{\textit{dept. name of organization (of Aff.)} \\
\textit{name of organization (of Aff.)}\\
City, Country \\
email address or ORCID}

}

\maketitle

\begin{abstract}
Panel time series data often contain heterogeneous temporal responses across entities, but existing neural forecasting models are not designed for observation of heterogeneous lag structures across entities. We propose AC-GATE, an Adaptive-Conditioning Encoder, which instantiates \emph{Conditional Moderated Distributed Lag} (CMDL) by utilizing observable entity-level proxy to condition lag-weight distributions over historical observations, thereby making effective lags as structural outputs of the model rather than post-hoc explanations.

The evaluation base on a layered audit protocol that separates predictive calibration from lag discovery. A synthetic panel with known ground-truth lags is used for mechanism recovery testing, and two real-world country-level panels used for external audit and stress testing. The results show that AC-GATE can recover heterogeneous lag structure in synthetic, and generate non-degenerate, externally structured effective lags in real data. The framework provides an auditable approach for studying entity-conditioned lag heterogeneity, without claiming causal identification or universal forecasting superiority.
\end{abstract}

\begin{IEEEkeywords}
Panel time series, heterogeneous lag discovery, entity-conditioned modeling, interpretable neural networks, gated sequence modeling, deep learning.
\end{IEEEkeywords}

\section{Introduction}
Panel time series data are more and more used to study processes like economic productivity, energy transition and climate-related outcomes. These processes cover changes of countries, regions, enterprises or organizations over time. This type of data contains two intertwined structures: temporal dependence in each entity and cross-sectional heterogeneity across entities. Classic research and recent surveys point out that the merging of dynamic relationships across entities can be misleading and reduce explainability when response parameters differ across units \cite{pesaran1995pmg,thayasivam2025panelsurvey}. These challenges are significant when the research not only focuses on the predictive accuracy of the model but also focuses on different effective lags of different entities in response to historical information.

In many real-world panel data, heterogeneous lags have practical meaning. In productivity analysis, inputs related to technology or ability may only impact outcomes after an adjustment period, organizational learning or complementary investments \cite{schweikl2020lessons}. In research on energy and emissions, the impacts of renewable energy, energy efficiency or institutional conditions may differ by country, development stage and emissions regime \cite{akram2020energy,mirziyoyeva2022re}. These observations indicate that single global lag structures may be too restrictive.

Current neural forecasting models provide flexible tools for nonlinear temporal prediction, but they are not designed for entity-conditioned lag discovery like Temporal Fusion Transformer (TFT) \cite{lim2021tft} and Gated Attention Network (GA-Net) \cite{xue2020gated}. As a result, the core question we consider in panel time series is: can the model discover systematically different effective lag structure across entities, and can the learned lag variation be validated by external entity-level features?

We propose AC-GATE, an Adaptive-Conditioning (AC) Encoder with a Scale-Invariant Lag Gate, which instantiates \emph{Conditional Moderated Distributed Lag} (CMDL) to discover entity-conditioned heterogeneous lags in panel time series. This structure uses observable entity-level proxy information to condition (and hence moderate) the past observations, thereby constructing weight distributions over lags based on entity characteristics. The result is an effective lag that is not a post-hoc explanation to a black-box, which is a structural output from the lag gate. This design is inspired by the absorptive-capacity view \cite{cohen1990absorptive}, while keeping agnostic to causal identification.

There is a layered audit protocol to evaluate AC-GATE, which separates predictive calibration from heterogeneous lag discovery. Meanwhile, the empirical evaluation covers three domains: synthetic, economic and energy, which have different evidential roles. In these domain, AC-GATE is compared against neural and econometric baselines, which include structural ablations that remove the AC encoder or disable the lag gate.

The contribution of our work has three aspects, First, formulating entity-conditioned heterogeneous lag discovery as a testable panel time-series mining problem, which is independence from standard prediction. Second, introducing AC-GATE, which generates entity-level effective lags via an explicit lag gating structure. Third, proposing a layered audit protocol that evaluates forecast calibration, lag non-degeneracy, external stratifier alignment and directional mechanism recovery. Results show that AC-GATE can recover true heterogeneous lag structure in synthetic panels, and identify structured lag heterogeneity in two real-world country-level panels.

\begin{figure}[!tbhp]
    \centering
    \includesvg[width=\columnwidth]{img/intro_use_case_box.svg}
    \caption{use case.}
    \label{fig:intro_case}
\end{figure}

\section{Related Work}\label{related_work}
Panel time series data contain temporal dependence and cross-sectional heterogeneity, so they cannot be considered as independent and identically distributed (i.i.d.) samples or ordinary uni-variate time series. The econometric literature shows that heterogeneity can display as discrete regimes or differences in distribution location or variable selection structure \cite{hansen1999thr,canay2011sim,babii2021mlpanel}. Recent review \cite{thayasivam2025panelsurvey} also emphasizes that the core challenges of panel time-series forecasting include cross-sectional heterogeneity, temporal dependence, non-stationary and explainability. These works indicate that the lag relationships in panel data should not be considered as homogeneous across entities, which should allow for entity-conditioned dynamic response differences. Therefore, even if existing panel forecasting frameworks can deal with heterogeneity, they usually do not produce verifiable entity-specific lag structures, which is the gap that this paper aims to fill.

Deep time-series models are good at capturing nonlinear dynamics but often lack explainability. The Temporal Fusion Transformer (TFT) combines static covariates, recurrent processing, attention and gating components for multi-horizon forecasting, which provides interpretability outputs \cite{lim2021tft}. On the other hand, the Gated Attention Network (GA-Net) shows that models can use input-dependent gates to select partially key elements, which can reduce redundant attention and improve interpretability \cite{xue2020gated}. And then, existing neural panel models have tried to put interpretable structure into the network, such as persistent change filters or nonlinear panel components, which can reveal auditable structure beyond forecasting \cite{yang2020panelnn,chronopoulos2023dnnpanel}. However, the TFT focuses on improving predictive performance and global explanation; GA-Net cannot produce a stable, cross-entity comparable lag-weight distribution.In general, current predictive models provide clues about explanation, rather than a modeling and validation framework for entity conditionation and using lag distribution as a structured output, so the lag gate + audit approach in this paper is designed for this purpose.

 
 
Except for architecture design, the application of panel machine learning also brings an evaluation problem. Supervised learning metrics mainly used to evaluate reserved predictive error, but it often fails to capture deployment-relevant objectives or the way of using information from the model when prediction alone is not the only goal \cite{lipton2018mythos}. For models that report explanatory structure, this issue is very important, because post-hoc explanations of black-box predictors can be incomplete or unfaithful, and interpretable models should exhibit internal structure that can be directly inspected and audited \cite{rudin2019stop}. In panel data, the temporal and cross-sectional dependencies further aggravate the complexity of the evaluation problem. Contemporaneous covariates, random unit-time splits, or preprocessing choices that use information unavailable at prediction time can lead to data leakage and exaggerate out-of-sample performance estimates \cite{cerqua2025mlpanel}. As a result, it is not enough to claim that a model has learned meaningful heterogeneous lag structure only based on high predictive accuracy. Our work aims to fill this evaluation gap by treating the learned effective lag as an auditable model-derived summary, rather than a post-hoc explanation or causal delay estimate.




\section{Problem Formulation}
% 这一段是面板设置
% This is panel setup, about concept abstract
\subsection{Panel Setup}
Considering a balanced panel time series with $N$ entities and $T$ time steps. For each entity $i \in \{1,\ldots,N\}$ and time $t \in \{1,\ldots,T\}$, a sequential covariate vector $X_{i,t}\in\mathbb{R}^{d_x}$ and a scalar target $Y_{i,t}\in\mathbb{R}$ can be observed. Each entity is associated with a time-invariant proxy vector $p_i\in\mathbb{R}^{M}$ and a static feature vector $s_i\in\mathbb{R}^{d_s}$ when available. This setting follows the standard panel time-series view where temporal dependence and cross-sectional heterogeneity are both central properties of the data \cite{thayasivam2025panelsurvey}. The proxy-as-conditioning view inspire by absorptive capacity, a kind of capacity to recognize, assimilate and apply external knowledge, which depends on prior related knowledge \cite{cohen1990absorptive}. Its core idea can be abstracted to a more general entity level (e.g., countries, regions or organizations): entities with different knowledge foundations, institutional conditions or development levels may absorb shocks or inputs with different speeds, and exhibit different delays in future outcomes.


% 对相关指标进行定义
% Define the relevant indicators
\subsection{Entity-Conditioned Heterogeneous Lag Mining}
Given a maximum lag horizon $K$, the goal is to learn an entity-specific lag-weight distribution:
\[
    \omega_i = (\omega_{i,1}, \ldots, \omega_{i,K}),
    \qquad
    \sum_{k=1}^{K}\omega_{i,k}=1,\quad \omega_{i,k}\ge 0.
\]
Where $\omega_{i,k}$ denotes the normalized weight assigned to lag $k$ for entity $i$.

The lag weights aggregate past sequential inputs into a per-entity context vector:
\begin{equation}
    c_{i,t} = \sum_{k=1}^{K}\omega_{i,k}X_{i,t-k}.
\end{equation}
And then, LSTM backbone processes the lag-context stream together with learned entity effects, static features, and optional macro controls to predict $\widehat{Y}_{i,t+1}$.

The interpretable byproduct is the per-entity effective lag:
\begin{equation}
    k_i^* = \sum_{k=1}^{K} k \, \omega_{i,k}.
\end{equation}
Because $\omega_i$ only depends on the entity-level conditioning score $z_i$, $k_i^*$ is constant across time for a fixed entity and can be audited. %可能也要删减

Separating heterogeneous lag mining from standard prediction. Forecasting evaluation focuses on accuracy of $\widehat{Y}_{i,t+1}$, yet lag mining evaluation focus on non-degenerate of the learned $k_i^*$ across entities and externally structured. The learned effective lag $k_i^*$ is evaluated through the audit protocol described in Section~\ref{sec:audit_protocol}.


% 讨论具体方法
% Discuss specific methods
\section{Method: AC-gate}
We propose AC-GATE, an Adaptive-Conditioning Encoder with a Scale-Invariant Lag Gate on top of a recurrent backbone. The model follows the pipeline:
\[
    p_i \rightarrow z_i \rightarrow \omega_i \rightarrow c_{i,t} \rightarrow \widehat{Y}_{i,t+1}.
\]
This design combines two reasons. First, the absorptive-capacity view we abstracted in section~\ref{related_work} suggests that entity-level prior knowledge can moderate how information is absorbed over time \cite{cohen1990absorptive}. Second, the sequence models base on gate and attention show that input-dependent selection can improve explainability by exposing which parts of a sequence or feature set are used by the model \cite{xue2020gated,lim2021tft}.

% AC architecture overview
\begin{figure}[!tbhp]
    \centering
    \includesvg[width=\columnwidth]{img/ac_gate_architecture.svg}
    \caption{Overview of the AC-GATE architecture.}
    \label{fig:ac_gate_arch}
\end{figure}


\subsection{Adaptive Conditioning Encoder}
For each entity $i$, the adaptive-conditioning encoder maps the entity-level proxy vector $p_i$ to a scalar absorption score:
\begin{equation}
    z_i = f_{\phi}(p_i),
    \qquad z_i\in\mathbb{R}.
\end{equation}
The proxy vector $p_i$ contains static entity-level covariates that capture persistent, time-invariant heterogeneity across entities. The scalar score $z_i$ is then used to condition the lag gate and is also projected into the initial recurrent state of the LSTM backbone.

A reconstruction head maps the scalar score back to the proxy space:
\begin{equation}
    \widehat{p}_i = g_{\psi}(z_i),
    \qquad \widehat{p}_i\in\mathbb{R}^{M}.
\end{equation}
Correspondingly, the reconstruction term can be used as an auxiliary agent reconstruction objective. It can as a lightweight diagnostic head or an explicit regularizer. In the detached variant, the reconstruction loss only trains the proxy decoder and should not be interpreted as directly regularizing the encoder representation. In the non-detached variant, the reconstruction loss encourages $z_i$ to retain information about the original proxy vector.


\subsection{Scale-Invariant Lag Gate}
Conditioned on the scalar absorption score $z_i$, the scale-invariant lag gate maps $z_i$ to an entity-specific lag-weight distribution:
\begin{equation}
\begin{aligned}
\omega_{i,k}
&= \mathrm{Softmax}_{k}\!\left(
\frac{g_{\theta}(z_i)_k - \lambda\, k/K}{\tau}
\right),\\
&\qquad k=1,\ldots,K.
\end{aligned}
\end{equation}
Where $g_{\theta}(z_i)\in\mathbb{R}^{K}$ outputs entity-specific lag logits, with $g_{\theta}(z_i)_k$ denoting the logit for lag $k$. \textit{Scalar-conditioned} means the logits are produced conditional on the scalar $z_i$. The operator $\mathrm{Softmax}_k(\cdot)$ normalizes these logits over $k=1,\ldots,K$ for each entity $i$.

The relative-position term $\lambda\cdot k/K$ penalizes farther lags on a normalized scale. The normalization by $K$ makes the position bias scale-invariant with respect to the maximum lag horizon. The temperature $\tau$ controls the sharpness of the resulting lag distribution.

\subsection{Backbone and Information-Flow Discipline}
The lag-weighted context $c_{i,t}$ is connected with a learned entity embedding, the optional static feature vector $s_i$, and optional macro-level controls. The sequence is processed by a two-layer LSTM, and the resulting hidden state is mapped to $\widehat{Y}_{i,t+1}$.

A key information-flow constraint is that the current sequential observation $X_{i,t}$ is not directly concatenated into the backbone input. Only lagged observations enter through $c_{i,t}$. This prevents a shortcut path where the recurrent backbone could ignore the lag gate and rely on contemporaneous covariates, which is a failure mode closely related to leakage concerns in panel machine learning \cite{cerqua2025mlpanel}.


% 改到这里
\subsection{Training Objective}
Training minimizes the domain-agnostic objective:
\begin{equation}
    \mathcal{L}
    =
    \mathcal{L}_{\mathrm{task}}
    +
    \lambda_r \mathcal{L}_{\mathrm{recon}}.
\end{equation}
The task loss $\mathcal{L}_{\mathrm{task}}$ is the forecasting mean squared error (MSE) between the predicted target $\widehat{Y}_{i,t+1}$ and the true target at the corresponding valid post-warmup step. The reconstruction loss $\mathcal{L}_{\mathrm{recon}}$ compares the original proxy vector $p_i$ with the reconstructed proxy vector $\widehat{p}_i$. The coefficient $\lambda_r$ controls the weight of the reconstruction term. 

\section{LAG DISCOVERY AND AUDIT PROTOCOL}\label{sec:audit_protocol}
A layered audit protocol is used to evaluate AC-GATE. This reflects the view that the evaluation of predictive models not only pay attention to test-set accuracy but also to whether they yield useful, deployment-relevant understanding (e.g., transparency, accountability, or faithful structure) beyond what standard predictive metrics capture \cite{lipton2018mythos,rudin2019stop}. This separation is necessary because forecast accuracy and structured lag discovery are related but different: a model maybe weak in a real-world data prediction while still producing a non-degenerate and externally structured effective-lag signal. Conversely, a strong forecasting baseline does not automatically imply an interpretable, entity-conditioned lag distribution. Algorithm~\ref{alg:acgate_audit} summarizes the four-layer audit protocol used throughout the experiments.

\begin{algorithm}[!tbhp]
\caption{Lag Discovery and Layered Audit}
\label{alg:acgate_audit}
\Input{$\{X_{i,t},Y_{i,t}\}$; $\{p_i\}$; optional $\{s_i\}$; lag horizon $K$; seeds $\mathcal{S}$; stratifiers $\Xi$; permutations $B$; threshold $\epsilon$; optional truth $k^{\mathrm{true}}_i$}
\Output{Forecast metrics; $\omega_i$; $k_i^\star$; L0--L3 verdicts}

\ForEach{$r\in\mathcal{S}$}{
Split chronologically\;
Fit normalizers on train only\;
Train AC-GATE by minimizing $\mathcal{L}_{\mathrm{task}}+\lambda_r\mathcal{L}_{\mathrm{recon}}$\;
Select checkpoint on validation\;

\tcp{L0 (forecast):}
Report test MSE/MAE/$R^2$\;

Compute effective lag $k_i^\star\gets\sum_{k=1}^{K}k\,\omega_{i,k}$\;

\tcp{L1 (degeneracy guard):}
Compute $\mathrm{sd}_i(k_i^\star)$\;
\If{$\mathrm{sd}_i(k_i^\star)\leq\epsilon$}{Mark seed degenerate\;\Continue}

\tcp{L2 (structured heterogeneity):}
\ForEach{$\xi^{(j)}\in\Xi$}{
Compute $\rho_{r,j}\gets\mathrm{Spearman}(k_i^\star,\xi_i^{(j)})$\;
Draw $B$ entity-label permutations of $\xi^{(j)}$ and compute $\rho^{(b)}$\;
Compute $p_{r,j}\gets(1+\#\{b:|\rho^{(b)}|\geq|\rho_{r,j}|\})/(B+1)$\;
}

\tcp{L3 (optional ground truth):}
\eIf{$k^{\mathrm{true}}_i$ available}{Report $\mathrm{MAE}(k^\star,k^{\mathrm{true}})$ and $\mathrm{Spearman}(k^\star,k^{\mathrm{true}})$\;}{Mark n/a\;}
}

Aggregate L0 across seeds\;
Aggregate L2 via $\mathrm{mean}(|\rho|)$, rejection share, and Fisher-combined $p$\;
Run paired seed-level tests vs. baselines/ablations\;
Assemble the L0--L3 verdict matrix\;
\end{algorithm}


\subsection{Audit Layers and Verdict Rules}
L0 focuses on the predictive calibration of each method. This layer is included to show whether the learned representation remains predicatively calibrated, but it is not the main mechanism-claim layer.

L1 is a degeneracy guard for lag discovery. For each seed, we compute the cross-entity standard deviation of the learned effective lag, $\mathrm{sd}_i(k_i^\star)$. If this value is below a small threshold $\epsilon$, the seed is marked as degenerate and is not used for L2 stratification. This rule prevents constant-lag models (e.g., uniform-lag or collapsed ablations) from being explained as evidence of heterogeneous lag structure.

L2 evaluates whether the learned effective lag is externally structured. For each pre-specified entity stratifier $\xi_i$, computing the seed-level Spearman correlation between $k_i^\star$ and $\xi_i$. And then, we compare its absolute value with an entity-label permutation. L2 evidence is summarized by mean $|\rho|$ across seeds, the fraction of seeds with permutation $p<0.05$, and a Fisher-combined
permutation $p$-value, which are shown in Table~\ref{tab:l2_sign_robust_quantities}. Stratifiers are only computed from the training window, so the test evaluates alignment with pre-existing entity characteristics rather than future information.


\begin{table}[!tbhp]
\centering
\caption{Sign-robust L2 summary statistics.}
\label{tab:l2_sign_robust_quantities}
\setlength{\tabcolsep}{5pt}
\renewcommand{\arraystretch}{1.1}
\begin{tabularx}{\columnwidth}{@{}l l@{}}
\toprule
Quantity & Evidence role \\
\midrule
$\mathrm{mean}(|\rho|)$ & Effect-size magnitude across random seeds \\
$\Pr\bigl(p_{\mathrm{perm}}<0.05\bigr)$ & Seed-level stability across random seeds \\
$p_{\mathrm{Fisher}}$ & Aggregated evidence across random seeds \\
\bottomrule
\end{tabularx}
\par\smallskip
{\footnotesize\raggedright\textit{Notes.} L2 uses absolute Spearman alignment because the sign of a learned latent lag score can vary across seeds. Permutation tests shuffle entity labels of $\xi_i$ while holding $k_i^\star$ fixed.\par}
\end{table}


L3 is only evaluated when ground-truth mechanism information is available. In the synthetic, we compare $k_i^\star$ with the true lag-generating target via MAE and Spearman correlation. In the real-world panels, L3 is marked as not applicable because no supervised mechanism label is observed. As a result, the learned lags are explained as model-derived effective-lag summaries, not as causal delay estimates.

\section{Experiments}
\subsection{Experimental Setup}


AC-GATE is evaluated on three domains with distinct evidential roles:
% owid_energy 可能需要同步进 refs.bib
\begin{itemize}
\item \textbf{Synthetic:}
Providing ground-truth heterogeneous lags, which is used for mechanism recovery. In the synthetic, script generates balanced panel time series under both linear and nonlinear lag-generating regimes. 
\item \textbf{PWT 11.0:}
 Economic panel data with the target variable $ctfp$ and an effective-labor-aware feature bundle, where 1980--2007 for training, 2008--2013 for validation, and 2014--2023 for testing. PWT provides internationally comparable measures of output, capital, human capital, labor input and productivity, which make it suitable for cross-country productivity analysis \cite{feenstra2015pwt}.
\item \textbf{OWID Energy \& WGI:}
Combining as energy data with the target is \texttt{co2\_per\_unit\_energy}, where 1996--2011 for training, 2012--2017 for validation, and 2018--2023 for testing. The WGI data provide governance-related entity stratifiers across six broad governance dimensions \cite{kaufmann2011wgi}.
\end{itemize}

To economic domain, prior work emphasizes cross-sectional parameter heterogeneity. Regime-dependent relationships, motivating designs that avoid relying solely on pooled mean effects \cite{canay2011sim,hansen1999thr,pesaran1995pmg}. And in the energy literature, the relationship between energy transition and emissions is also documented to vary across countries and quantiles \cite{akram2020energy,mirziyoyeva2022re}.



To support reproducibility and ensure fair comparisons, we summarize the baseline/ablation design choices that define the controlled comparison in Table~\ref{tab:fair_baselines}. Baselines include Plain LSTM, TFT \cite{lim2021tft}, GA-Net \cite{xue2020gated}, and Grouped ARDL\cite{pesaran1995pmg,babii2021mlpanel}. Structural ablations are No-AC-Encoder ($z_i$ fixed), Uniform-Lag ($\omega_{i,k}=1/K$), and No-Recon-Reg ($\lambda_r=0$). Details of experimental parameters and data preprocessing can be found in the appendix~\ref{app:repro}. In addition, negative-control check via proxy permutation also is introduced to diagnose spurious predictability, which is shown in Table~\ref{tab:proxy_shuffle_setup}


\begin{table}[!tbhp]
\caption{Fair-comparison design.}
\label{tab:fair_baselines}
\centering
\scriptsize
\setlength{\tabcolsep}{3pt}
\renewcommand{\arraystretch}{1.03}
\begin{tabular}{p{0.22\columnwidth} p{0.72\columnwidth}}
\toprule
Method & Difference (all share the same splits, $K$, and protocol) \\
\midrule
AC-GATE & Full model (AC encoder + lag gate + recon.). \\
Plain LSTM & No AC encoder; no lag gate (matched backbone). \\
TFT & Transformer with static covariates and temporal attention. \\
GA-Net & Gated-attention baseline. \\
No-AC-Encoder & Fix $z_i$ (no entity-conditioned lag). \\
Uniform-Lag & Uniform $\omega_{i,k}$ (no learned lag heterogeneity). \\
No-Recon-Reg & Set $\lambda_r=0$ (no recon.). \\
Grouped ARDL & Econ./Energy only: grouped OLS distributed-lag baseline. \\
\bottomrule
\end{tabular}
\par\smallskip
{\footnotesize\raggedright\textit{Notes.} Plain LSTM as same as AC-GATE backbone capacity, but removes the adaptive-conditioning encoder and lag gate. Grouped ARDL uses the same feature sets and the same $K$-lag block, and grouping only uses training-window information.\par}
\end{table}

\begin{table}[!tbhp]
\centering
\caption{Proxy-Shuffle Negative Control Setup}
\label{tab:proxy_shuffle_setup}
\footnotesize
\setlength{\tabcolsep}{3pt}
\renewcommand{\arraystretch}{1.05}
\begin{tabular}{@{}l l@{}}
\toprule
Aspect & Setting \\
\midrule
Purpose & Test if L2 alignment needs correct proxy--entity match. \\
Control type & Information negative control (not ablation). \\
Shuffled component & Permute entity proxy vectors across entities. \\
Preserved components & Data/IDs, splits/strata, arch., and hparams fixed. \\
Permutation rule & Per seed: one derangement (no self-proxy). \\
Training protocol & Retrain AC-GATE for 20 seeds (same protocol). \\
Evaluation protocol & Recompute L2 stratified $k^*$ audit (orig. stratifiers). \\
Primary metric & Mean $|\rho|$ drop; reject share; Fisher combined p. \\
Interpretation & Collapse implies alignment is proxy-semantic. \\
\bottomrule
\end{tabular}
\end{table}


\subsection{Paired significance tests}
All models are evaluated via 20 random seeds. For paired model comparisons, we use a paired Wilcoxon test on seed-level metrics, which compares AC-GATE against each baseline or ablation model under the same seeds. This scheme detects the systematic deviation of paired differences from zero while controlling for training variability across seeds.



For the synthetic experiments, the main test metric is $k^\star$ MAE because the ground-truth effective lag is known. For the real-data experiments, we do not use Wilcoxon tests to validate real-data mechanisms, paired tests are applied only to forecast-layer metrics like $R^2$. Stratifier alignment is evaluated separately through entity-level permutation tests in each seed and Fisher aggregation across seeds. Detailed paired-test results are reported in Appendix~\ref{app:wilcoxon}.



\subsection{Experimental results and analysis}

\begin{figure}[!!tbhp]
    \centering    % In a two-column layout, \textwidth spans both columns; use \columnwidth for a single column figure.
    \includegraphics[width=\columnwidth]{img/realdata_case_study_panel.png}
    \caption{Relationship between entity-level stratifiers and learned effective lag.}
    {\footnotesize\raggedright\textit{Notes.} Only use training window information, higher levels of human capital/rule of law correspond to shorter learned effective lags. This indicates that the national-level lag structure learned by AC-GATE is not random, which exhibits a systematic stratification related to the development/governance status during the training period.\par}
    \label{fig:case_study_panel}
\end{figure}

\begin{figure*}[!t]
    \centering
    \includesvg[width=\textwidth]{img/compact_evidence_panel.svg}
    \caption{Evidence for AC-GATE recovery, forecast-mechanism decoupling, and stratifier alignment.}
    \label{fig:compact_evidence_panel}
    {\footnotesize\raggedright\textit{Notes.} Panel A summarizes synthetic mechanism recovery under linear and nonlinear settings. Panel B contrasts predictive performance with the strength of learned stratifier association in real-data tasks. Panel C reports AC-GATE stratifier alignment across economics and energy variables, with the right column showing the percentage of seeds whose permutation-test p-value is below 0.05.\par}
\end{figure*}

In Fig~\ref{fig:case_study_panel}, AC-GATE's learned entity-level effective lag \(k_i^\star\) is negatively associated with external stratifiers: human capital in the economics panel (Spearman \(\rho=-0.63\), \(p\approx 3\times 10^{-11}\), \(n=88\)) and rule of law in the energy panel (\(\rho=-0.61\), \(p\approx 3\times 10^{-9}\), \(n=78\)). This consistency suggests that more stronger-institution entities are assigned shorter effective memory, which support stable domain-semantic alignment of the recovered heterogeneous lag.

The Fig.~\ref{fig:compact_evidence_panel} is the evidence for AC-GATE recovery. The synthetic experiments test recovery of known heterogeneous lag structure. AC-GATE achieves strong rank alignment between learned and true effective lags (Spearman $\rho \approx 0.945$ linear; $0.907$ nonlinear), while No-AC-Encoder and Uniform-Lag show near-zero alignment due to structurally collapsed entity-level lag variation.

In the economics domain, the learned lag mechanism is non-degenerate and stratified. AC-GATE maintains entity-level lag variation ($k^\star$ std. $\approx 0.167$) and lag-gate sensitivity, while No-AC-Encoder and Uniform-Lag collapse the structured mechanism. Moreover, learned effective lags of AC-GATE align with economically meaningful stratifiers: human capital shows the strongest association (mean $|\rho| \approx 0.371$), followed by GDP per worker ($0.278$) and capital per worker ($0.257$), with 65--80\% of seeds significant under the permutation test.

In the energy domain, the mechanism diagnostics has reference value: AC-GATE and No-Recon keep non-degenerate lag variation, whereas No-AC-Encoder and Uniform-Lag flatten the entity-level lag structure again. Learned effective lags of AC-GATE show strong alignment with governance and development stratifiers, with mean $|\rho| \approx 0.735$ for rule of law, $0.716$ for government effectiveness, and $0.609$ for GDP per capita. These associations are stable across seeds, with 90--95\% of seeds passing the permutation test.

Fig~\ref{fig:r2_floor} shows that the forecast \(R^2\) values in real-data panels cluster around a narrow band near zero for AC-GATE and all deep learning baselines, except for the Grouped ARDL in the energy domain which achieves a significantly higher \(R^2\). This represents the empirical \(R^2\) close to zero is a characteristic of the forecasting itself, which motivates the separation of forecast calibration from structured lag mechanism evidence.

\begin{figure}[!tbhp]
  \centering
  \includesvg[width=\columnwidth]{img/panel_r2_floor.svg}
  \caption{Forecast $R^2$ on the two real-world panels.}
  {\footnotesize\raggedright\textit{Notes.} All deep baselines (incl. AC-GATE and ablations) fall in the near-zero band $[-0.05, 0.10]$. Grouped ARDL flips sign across domains ($-0.089$ in Economics; $+0.607$ in Energy), so forecast $R^2$ alone gives unstable rankings; use mechanism evidence (Sec.~\ref{sec:audit_protocol}, L1/L2).\par}
  \label{fig:r2_floor}
\end{figure}

As shown in Table~\ref{tab:proxy_shuffle_compact}, the proxy-shuffle negative control destroys the semantic correspondence between entities and their proxy vectors in real-data panels. In this control, L2 alignment fails and predictive performance remains nearly unchanged. This indicates that the stratified \(k^\star\) in real domains is driven by correct entity-proxy conditioning, not by generic model capacity or forecasting performance.
\begin{table}[!tbhp]
\centering
\caption{Proxy-shuffle negative control on real-data L2 alignment.}
\label{tab:proxy_shuffle_compact}
% Avoid unconditional down-scaling (which makes fonts tiny):
% set the table font explicitly and tighten column spacing instead.
\scriptsize
\setlength{\tabcolsep}{2pt}
\renewcommand{\arraystretch}{1.05}
\begin{tabular}{llccccc}
\toprule
Dom. & Model & Mean $|\rho|$ & $\mathrm{sd}(k^\star)$ & Frac. $p<5\%$ & Fisher $p$ & Test $R^2$ \\
\midrule
Econ. & AC-GATE & 0.302 & 0.168 & 0.72 & $10^{-46}$--$10^{-24}$ & 0.054 \\
Econ. & Proxy-shuf. & 0.058 & 0.184 & 0.00 & 0.94--1.00 & 0.055 \\
Ener. & AC-GATE & 0.686 & 0.109 & 0.92 & $10^{-79}$--$10^{-77}$ & -0.029 \\
Ener. & Proxy-shuf. & 0.088 & 0.073 & 0.07 & 0.39--0.75 & -0.028 \\
\bottomrule
\end{tabular}
\vspace{1mm}
\begin{flushleft}
\footnotesize \textit{Note}: Economics stratifiers are human capital, capital per worker, and GDP per worker. Energy stratifiers are government effectiveness, GDP per capita, and rule of law. Proxy shuffling preserves model architecture and training protocol but breaks the entity--proxy correspondence.
\end{flushleft}
\end{table}


Taken together, the experiments in real-data show that strongest evidence of AC-GATE lies in structured and non-degenerate mechanism recovery, rather than universal predicting superiority. Base on the experimental results and audit protocol, the verdict matrix shown in the Table~\ref{tab:verdict_matrix}.



\begin{table}[!tbhp]
\centering
\caption{Verdict matrix.}
\label{tab:verdict_matrix}
\renewcommand{\arraystretch}{1.25}
\setlength{\tabcolsep}{6pt}
\begin{tabularx}{\columnwidth}{l *{4}{>{\centering\arraybackslash}X}}
\toprule
Domain & L0 & L1 & L2 & L3 \\
\midrule
Synthetic & \nc & \yes & \yes & \yes \\
Economics & \nc & \yes & \yes & \NA \\
Energy & \nc & \yes & \yes & \NA \\
\bottomrule
\end{tabularx}
\par\smallskip
{\footnotesize\raggedright Green denotes supported, amber not certified, red ruled out, and gray not claimed.\par}
\end{table}




\section{Discussion}
Table~\ref{tab:verdict_matrix} summarizes the layered verdict, which shows that AC-GATE is supported on L1/L2 in all domains.
The experiments show that AC-GATE is not a universal forecasting superior. In the synthetic with known real heterogeneous lag structure, AC-GATE has high rank consistency in recovering the true effective lag. However, in real-data data, the predictive performance of AC-GATE does not show advantage. This distinction is crucial to the boundaries of AC-GATE, it is designed to reveal non-degenerate, entity-conditioned lag structure in panel data. Therefore, the real-data results support explanation of mechanism and not claim a universal forecasting superior. This separates the present work from causal panel learning methods that explicitly target heterogeneous treatment effects and counterfactual recovery \cite{zhou2025codeal}.

The learned $k^\star$ should be regarded as a model-derived effective lag summary, not as a causal delay estimate. In synthetic, $k^\star$ can be evaluated directly against the known data-generating lag. In real panels, the analogous question is still testable: whether the learned entity-level lags are non-randomly stratified by pre-specified external indicators. The economics results show that learned lags align with development-related indicators (e.g., human capital, GDP per worker and capital per worker), and the energy results show even stronger alignment with governance and development indicators (e.g., rule of law, government effectiveness and GDP per capita). These findings suggest that AC-GATE can capture structured heterogeneity in temporal response patterns, but they do not establish directional causality in economy or environment.

The ablation results clarify where the mechanism evidence comes from. Removing the AC encoder or replacing the lag gate with a uniform lag distribution eliminates entity-level lag variation, thereby making stratifier alignment impossible. By contrast, removing the reconstruction regularizer hardly changes synthetic lag recovery, which indicate that the reconstruction term is an auxiliary stabilizer, rather than the source of the recovered mechanism. As a result, the essential inductive bias is the combination of entity conditioning and explicit lag gating: the model must both represent entity-level context and convert that context into a non-uniform lag distribution.


\section{Conclusion \& Future Work}
Based on the current results, the primary contribution of AC-GATE is to identify and characterize non-degenerate, entity-conditioned lag heterogeneity in panel data, which shows good recovery of true effective lags in synthetic. It is important to emphasize that AC-GATE is not designed to improve predictive accuracy, its output $k^\star$ should be regarded as a  model-derived effective lag summary, not as a directly explainable causal delay estimate. As a result, AC-GATE is more suitable as a mechanism diagnostic and interpretability tool, which can provide auditable lag structure evidence on existing predictive models, while also can be combined with external stratifiers to provide clues for subsequent mechanism discussion and causal identification, rather than giving direct causal conclusions.


There are some aspects that can be considered for future work. First, further exploring the performance of AC-GATE on different types of panel data, and its applicability to different prediction tasks. Second, trying to introduce more complex entity conditioning encoders or more flexible lag gating mechanisms to enhance the ability of the model to capture complex lag structures. Finally, combining domain knowledge to design more targeted external stratifiers, and validating the relationship between the lag structure discovered by AC-GATE and actual mechanisms in more empirical backgrounds, to further enhance its explanatory power and practical value.



\appendices

\section{Paired Wilcoxon Significance Tests}\label{app:wilcoxon}
\subsection{Synthetic: \(k^\star\) MAE}
\begin{table}[!tbhp]
\centering
\caption{Paired Wilcoxon test of \(k^\star\) MAE on the synthetic panel
(reference = AC-GATE, 20 seeds).}
\label{tab:app_wilcoxon_synth_kstar}
\footnotesize
\setlength{\tabcolsep}{4pt}
\noindent\begin{tabular*}{\columnwidth}{@{\extracolsep{\fill}} l l r r r r}
\toprule
Scenario & Method & Mean diff. & Median diff. & $W$ & $p$ \\
\midrule
linear    & No AC Encoder       & $+0.772$ & $+0.785$ & $0$  & $1.9{\times}10^{-6}$ \\
linear    & No Recon Reg.       & $+0.000$ & $+0.000$ & $86$ & $0.498$ \\
linear    & Plain LSTM          & $+0.548$ & $+0.588$ & $0$  & $1.9{\times}10^{-6}$ \\
linear    & TFT                 & $+0.508$ & $+0.545$ & $0$  & $1.9{\times}10^{-6}$ \\
linear    & GA-Net              & $+0.515$ & $+0.526$ & $0$  & $1.9{\times}10^{-6}$ \\
linear    & Uniform Lag         & $+0.754$ & $+0.799$ & $0$  & $1.9{\times}10^{-6}$ \\
nonlinear & No AC Encoder       & $+0.968$ & $+0.972$ & $0$  & $1.9{\times}10^{-6}$ \\
nonlinear & No Recon Reg.       & $-0.000$ & $-0.000$ & $80$ & $0.368$ \\
nonlinear & Plain LSTM          & $+1.144$ & $+1.122$ & $0$  & $1.9{\times}10^{-6}$ \\
nonlinear & TFT                 & $+1.251$ & $+1.264$ & $0$  & $1.9{\times}10^{-6}$ \\
nonlinear & GA-Net              & $+1.222$ & $+1.226$ & $0$  & $1.9{\times}10^{-6}$ \\
nonlinear & Uniform Lag         & $+1.635$ & $+1.669$ & $0$  & $1.9{\times}10^{-6}$ \\
\bottomrule
\end{tabular*}
\par\smallskip
{\footnotesize\raggedright\textit{Notes.} Mean/median diff.\ $=$ method $-$ AC-GATE.\par}
\end{table}

\subsection{Synthetic: Forecast Task Loss}
% same structure as Table~\ref{tab:app_wilcoxon_synth_kstar},
% from synthetic_significance_task_loss.csv
\begin{table}[!tbhp]
\centering
\caption{Paired Wilcoxon test of forecast task loss on the synthetic
panel (reference = AC-GATE, 20 seeds; lower is better).}
\label{tab:app_wilcoxon_synth_taskloss}
\footnotesize
\setlength{\tabcolsep}{4pt}
\noindent\begin{tabular*}{\columnwidth}{@{\extracolsep{\fill}} l l r r r r}
\toprule
Scenario & Method & Mean diff. & Median diff. & $W$ & $p$ \\
\midrule
linear    & No AC Encoder       & $+0.0267$ & $+0.0267$ & $0$  & $1.9{\times}10^{-6}$ \\
linear    & No Recon Reg.       & $+0.0001$ & $-0.0000$ & $98$ & $0.812$ \\
linear    & Plain LSTM          & $+0.0333$ & $+0.0338$ & $0$  & $1.9{\times}10^{-6}$ \\
linear    & TFT                 & $-0.0043$ & $-0.0032$ & $15$ & $2.6{\times}10^{-4}$ \\
linear    & GA-Net              & $-0.0037$ & $-0.0065$ & $59$ & $0.090$ \\
linear    & Uniform Lag         & $+0.0333$ & $+0.0333$ & $0$  & $1.9{\times}10^{-6}$ \\
nonlinear & No AC Encoder       & $+0.0282$ & $+0.0352$ & $0$  & $1.9{\times}10^{-6}$ \\
nonlinear & No Recon Reg.       & $+0.0000$ & $-0.0000$ & $85$ & $0.475$ \\
nonlinear & Plain LSTM          & $+0.0549$ & $+0.0543$ & $0$  & $1.9{\times}10^{-6}$ \\
nonlinear & TFT                 & $-0.0022$ & $-0.0007$ & $83$ & $0.430$ \\
nonlinear & GA-Net              & $+0.0029$ & $+0.0044$ & $33$ & $5.6{\times}10^{-3}$ \\
nonlinear & Uniform Lag         & $+0.0560$ & $+0.0585$ & $0$  & $1.9{\times}10^{-6}$ \\
\bottomrule
\end{tabular*}
\par\smallskip
{\footnotesize\raggedright\textit{Notes.} Mean/median diff.\ $=$ method $-$ AC-GATE.\par}
%Source: \texttt{synthetic\_significance\_task\_loss.csv}.\par}
\end{table}



\subsection{Real Domains: Test \(R^2\)}
% Combine economics_significance_test_r2.csv and
% energy_significance_test_r2.csv into one or two tables.
\begin{table}[!tbhp]
\centering
\caption{Paired Wilcoxon test of test \(R^2\) on the real-data panels
(reference = AC-GATE, 20 seeds; higher is better).}
\label{tab:app_wilcoxon_real_r2}
\footnotesize
\setlength{\tabcolsep}{4pt}
\noindent\begin{tabular*}{\columnwidth}{@{\extracolsep{\fill}} l l r r r r}
\toprule
Domain & Method & Mean diff. & Median diff. & $W$ & $p$ \\
\midrule
Economics & Grouped ARDL  & $-0.1433$ & $-0.1447$ & $0$  & $1.9{\times}10^{-6}$ \\
Economics & No AC Encoder & $-0.0023$ & $-0.0007$ & $68$ & $0.177$ \\
Economics & No Recon Reg. & $+0.0011$ & $-0.0009$ & $78$ & $0.330$ \\
Economics & Plain LSTM    & $+0.0483$ & $+0.0413$ & $9$  & $6.3{\times}10^{-5}$ \\
Economics & TFT           & $-0.0155$ & $-0.0240$ & $64$ & $0.133$ \\
Economics & GA-Net        & $-0.0164$ & $-0.0258$ & $56$ & $0.070$ \\
Economics & Uniform Lag   & $+0.0498$ & $+0.0410$ & $0$  & $1.9{\times}10^{-6}$ \\
\midrule
Energy    & Grouped ARDL  & $+0.6356$ & $+0.6352$ & $0$  & $1.9{\times}10^{-6}$ \\
Energy    & No AC Encoder & $+0.0002$ & $+0.0001$ & $97$ & $0.784$ \\
Energy    & No Recon Reg. & $+0.0000$ & $+0.0000$ & $24$ & $1.4{\times}10^{-3}$ \\
Energy    & Plain LSTM    & $-0.0002$ & $-0.0020$ & $95$ & $0.729$ \\
Energy    & TFT           & $-0.0048$ & $-0.0046$ & $36$ & $8.3{\times}10^{-3}$ \\
Energy    & GA-Net        & $-0.0046$ & $-0.0055$ & $35$ & $7.3{\times}10^{-3}$ \\
Energy    & Uniform Lag   & $-0.0013$ & $+0.0002$ & $86$ & $0.498$ \\
\bottomrule
\end{tabular*}
\par\smallskip
{\footnotesize\raggedright\textit{Notes.} Mean/median diff.\ $=$ method $-$ AC-GATE.\par}
%Sources: \texttt{economics\_significance\_test\_r2.csv},
%\texttt{energy\_significance\_test\_r2.csv}.\par}
\end{table}



\section{Full Multi-seed and Audit Summaries}\label{app:audit_full}

\FloatBarrier
\subsection{Synthetic L0/L3 Summary}
\par\medskip
\begin{table}[!tbhp]
\caption{Synthetic 20-seed summary. Task loss is the forecasting MSE;
\(k^\star\) MAE and Spearman \(\rho(k^\star,k^\star_{\mathrm{true}})\)
are L3 mechanism-recovery indicators.}
\label{tab:app_synth_summary}
\scriptsize
\setlength{\tabcolsep}{3pt}
% Left aligned; make the tabular block exactly \columnwidth (consistent with other three-line tables)
\noindent\begin{tabular*}{\columnwidth}{@{\extracolsep{\fill}} l l r r r}
\toprule
Scenario & Method & Task loss & \(k^\star\) MAE & Spearman \(\rho\) \\
\midrule
linear & AC-GATE              & $0.0362\,(0.0044)$ & $1.159\,(0.227)$ & $\;\;0.945\,(0.078)$ \\
linear & No Recon Reg.        & $0.0363\,(0.0043)$ & $1.159\,(0.227)$ & $\;\;0.945\,(0.078)$ \\
linear & Plain LSTM           & $0.0695\,(0.0029)$ & $1.707\,(0.091)$ & $\;\;0.356\,(0.066)$ \\
linear & TFT                  & $0.0319\,(0.0025)$ & $1.668\,(0.087)$ & $\;\;0.472\,(0.047)$ \\
linear & GA-Net               & $0.0325\,(0.0069)$ & $1.674\,(0.087)$ & $\;\;0.508\,(0.065)$ \\
linear & No AC Encoder        & $0.0629\,(0.0054)$ & $1.931\,(0.162)$ & $\;\;0.000\,(0.000)$ \\
linear & Uniform Lag          & $0.0695\,(0.0027)$ & $1.913\,(0.088)$ & $\;\;0.000\,(0.000)$ \\
\midrule
nonlin.& AC-GATE              & $0.0378\,(0.0065)$ & $1.467\,(0.248)$ & $\;\;0.907\,(0.212)$ \\
nonlin.& No Recon Reg.        & $0.0378\,(0.0067)$ & $1.467\,(0.248)$ & $\;\;0.909\,(0.205)$ \\
nonlin.& Plain LSTM           & $0.0927\,(0.0061)$ & $2.611\,(0.138)$ & $\;\;0.344\,(0.092)$ \\
nonlin.& TFT                  & $0.0356\,(0.0047)$ & $2.718\,(0.094)$ & $\;\;0.474\,(0.095)$ \\
nonlin.& GA-Net               & $0.0408\,(0.0029)$ & $2.689\,(0.091)$ & $\;\;0.477\,(0.067)$ \\
nonlin.& No AC Encoder        & $0.0661\,(0.0184)$ & $2.435\,(0.261)$ & $\;\;0.000\,(0.000)$ \\
nonlin.& Uniform Lag          & $0.0938\,(0.0081)$ & $3.102\,(0.098)$ & $\;\;0.000\,(0.000)$ \\
\bottomrule
\end{tabular*}
\par\smallskip
{\footnotesize\raggedright\textit{Notes.} Cells report mean (sd) over 20 seeds.\par}
%Source: \texttt{synthetic\_multiseed\_summary.csv}.\par}
\end{table}

\FloatBarrier
\subsection{Real-Domain Compact Summaries (L0 \& L1)}
\par\medskip
\begin{table}[!tbhp]
\centering
\caption{Economics 20-seed compact summary. }
\label{tab:app_econ_compact}
\scriptsize
\setlength{\tabcolsep}{3pt}
%\begin{tabular}{l r r r r}
\noindent\begin{tabular*}{\columnwidth}{@{\extracolsep{\fill}} l r r r r}
\toprule
Method & Test \(R^2\) & \(k^\star_{\sigma}\) & \(\omega\)-entropy & Gate sensitivity \\
\midrule
AC-GATE        & $\;\;0.0541\,(0.0373)$ & $0.167$ & $1.853$ & $0.499$ \\
No Recon Reg.  & $\;\;0.0553\,(0.0329)$ & $0.109$ & $1.844$ & $0.339$ \\
No AC Encoder  & $\;\;0.0518\,(0.0354)$ & $0.000$ & $1.839$ & $0.000$ \\
Uniform Lag    & $\;\;0.1040\,(0.0280)$ & $0.000$ & $2.303$ & $0.735$ \\
Plain LSTM     & $\;\;0.1024\,(0.0214)$ & $1.894$ & $1.873$ & --- \\
TFT            & $\;\;0.0386\,(0.0251)$ & $1.666$ & $1.968$ & --- \\
GA-Net         & $\;\;0.0377\,(0.0133)$ & $1.769$ & $1.912$ & --- \\
Grouped ARDL   & $-0.0891\,(0.0000)$    & ---     & ---     & --- \\
\bottomrule
\end{tabular*}

\par\smallskip
{\footnotesize\raggedright\textit{Notes.} \(k^\star_{\sigma}\) is the
mean of the cross-entity standard deviation of \(k^\star\) per seed
(L1); a value close to zero indicates a degenerate lag distribution.\par}
\end{table}

\begin{table}[!tbhp]
\caption{Energy 20-seed compact summary (analogous to Table~\ref{tab:app_econ_compact}).}
\label{tab:app_energy_compact}
\scriptsize
\setlength{\tabcolsep}{3pt}
% Left aligned; make the tabular block exactly \columnwidth (consistent with other three-line tables)
\noindent\begin{tabular*}{\columnwidth}{@{\extracolsep{\fill}} l r r r r}
\toprule
Method & Test \(R^2\) & \(k^\star_{\sigma}\) & \(\omega\)-entropy & Gate sensitivity \\
\midrule
AC-GATE        & $-0.0286\,(0.0050)$ & $0.108$ & $1.857$ & $0.269$ \\
No Recon Reg.  & $-0.0286\,(0.0050)$ & $0.108$ & $1.857$ & $0.269$ \\
No AC Encoder  & $-0.0284\,(0.0046)$ & $0.000$ & $1.845$ & $0.000$ \\
Uniform Lag    & $-0.0299\,(0.0035)$ & $0.000$ & $2.303$ & $0.394$ \\
Plain LSTM     & $-0.0288\,(0.0042)$ & $1.331$ & $2.059$ & --- \\
TFT            & $-0.0334\,(0.0055)$ & $1.730$ & $1.775$ & --- \\
GA-Net         & $-0.0331\,(0.0042)$ & $1.419$ & $1.993$ & --- \\
Grouped ARDL   & $\;\;0.6071\,(0.0000)$ & ---  & ---     & --- \\
\bottomrule
\end{tabular*}
\par\smallskip
{\footnotesize\raggedright\textit{Notes.} Cells: mean (sd) over 20 seeds.\par}
%Sources: \texttt{economics\_compact\_summary.csv},
%\texttt{energy\_compact\_summary.csv}.\par}
\end{table}

\FloatBarrier
\subsection{L2 Stratified \(k^\star\) (All Stratifiers)}
\begin{table}[!tbhp]
\centering
\caption{Economics L2 stratified \(k^\star\) summary (only methods that pass L1
yield valid rows).}
\label{tab:app_econ_l2_full}
\scriptsize
\setlength{\tabcolsep}{3pt}
\noindent\begin{tabular*}{\columnwidth}{@{\extracolsep{\fill}} l l r r r r}
\toprule
Method & Stratifier & \(\overline{|\rho|}\) & median \(\rho\)
       & \(\Pr(p{<}.05)\) & \(p_{\mathrm{Fisher}}\) \\
\midrule
AC-GATE & HC                         & $0.371$ & $-0.296$ & $0.80$ & $1.0{\times}10^{-46}$ \\
AC-GATE & log GDP per worker         & $0.278$ & $-0.146$ & $0.70$ & $3.5{\times}10^{-37}$ \\
AC-GATE & log capital per worker     & $0.257$ & $-0.135$ & $0.65$ & $1.4{\times}10^{-24}$ \\
No Recon Reg. & HC                   & $0.418$ & $-0.298$ & $0.80$ & $1.5{\times}10^{-50}$ \\
No Recon Reg. & log GDP per worker   & $0.287$ & $-0.126$ & $0.65$ & $2.5{\times}10^{-31}$ \\
No Recon Reg. & log capital per worker & $0.272$ & $-0.124$ & $0.60$ & $6.1{\times}10^{-30}$ \\
\bottomrule
\end{tabular*}
\par\smallskip
{\footnotesize\raggedright\textit{Notes.} No-AC-Encoder, Uniform-Lag, Plain LSTM, TFT, and GA-Net are L1-degenerate (all 20 seeds invalid) and are omitted.\par}
%Source: \texttt{economics\_stratified\_kstar\_aggregated.csv}.\par}
\end{table}

\begin{table}[!tbhp]
\centering
\caption{Energy L2 stratified \(k^\star\) summary.}
\label{tab:app_energy_l2_full}
\scriptsize
\setlength{\tabcolsep}{3pt}
\noindent\begin{tabular*}{\columnwidth}{@{\extracolsep{\fill}} l l r r r r}
\toprule
Method & Stratifier & \(\overline{|\rho|}\) & median \(\rho\)
       & \(\Pr(p{<}.05)\) & \(p_{\mathrm{Fisher}}\) \\
\midrule
AC-GATE & rule of law              & $0.735$ & $\;\;0.237$ & $0.95$ & $1.8{\times}10^{-79}$ \\
AC-GATE & gov.\ effectiveness      & $0.716$ & $\;\;0.129$ & $0.90$ & $9.0{\times}10^{-77}$ \\
AC-GATE & log GDP per capita       & $0.609$ & $\;\;0.185$ & $0.90$ & $1.5{\times}10^{-77}$ \\
No Recon Reg. & rule of law        & $0.735$ & $\;\;0.237$ & $0.95$ & $1.8{\times}10^{-79}$ \\
No Recon Reg. & gov.\ effectiveness& $0.716$ & $\;\;0.129$ & $0.90$ & $9.0{\times}10^{-77}$ \\
No Recon Reg. & log GDP per capita & $0.609$ & $\;\;0.185$ & $0.90$ & $1.5{\times}10^{-77}$ \\
\bottomrule
\end{tabular*}
\par\smallskip
{\footnotesize\raggedright\textit{Notes.} No-AC-Encoder, Uniform-Lag, Plain LSTM, TFT, and GA-Net are L1-degenerate (all 20 seeds invalid) and are omitted. \(\overline{|\rho|}\) is sign-robust; the near-zero seed-mean of signed \(\rho\) reflects the seed-wise sign-flip invariance discussed in Section~\ref{sec:audit_protocol} (L3).\par}
%Source: \texttt{energy\_stratified\_kstar\_aggregated.csv}.\par}
\end{table}

\FloatBarrier
\section{Reproducibility Metadata}\label{app:repro}

\FloatBarrier
\subsection{Environment}
\begin{table}[!tbhp]
\centering
\caption{Experiment environment for the locked 20-seed suite.}
\label{tab:app_env}
\footnotesize
\setlength{\tabcolsep}{4pt}
\renewcommand{\arraystretch}{1.05}
\noindent\begin{tabular*}{\columnwidth}{@{\extracolsep{\fill}} l l}
\toprule
Item & Value \\
\midrule
OS                       & Windows 11 (10.0.26200) \\
Python                   & 3.13.9 \\
PyTorch                  & 2.9.1+cu130 \\
CUDA available           & yes \\
Device                   & \texttt{cuda} (NVIDIA GeForce RTX 5080) \\
\texttt{cudnn.deterministic} & false \\
\texttt{cudnn.benchmark}     & false \\
Deterministic algorithms & off \\
%Conda env spec           & \texttt{environment.yml} \\
%Pip lock                 & \texttt{requirements.txt} \\
\bottomrule
\end{tabular*}
\end{table}


\FloatBarrier
\section{Reproducibility Metadata}\label{app:repro}

\FloatBarrier
\subsection{Experiment Parameters}
\begin{table}[!tbhp]
\centering
\caption{Locked experiment parameters.}
\label{tab:app_experiment_parameters}
\setlength{\tabcolsep}{4pt}
\renewcommand{\arraystretch}{1.05}
\noindent\begin{tabular*}{\columnwidth}{@{\extracolsep{\fill}} l l l}
\toprule
Scope & Block & Value \\
\midrule
Shared & Seeds & 20 (0--19) \\
Shared & Core & \(K=10\), \(d=64\), LSTM=2 \\
Shared & Optim. & drop 0.05; lr \(10^{-3}\); pat 20 \\
Shared & Lag gate & \(\lambda_r=0.1\); \(\tau=1\); softmax; b=1 \\
Shared & Runtime & clip 1; device auto \\
Syn. & Plan & lin+nonlin; ep 200; val 0.2 \\
Econ. & Contract & target ctfp; eff.-labor bundle \\
Econ. & Reg. & anchor recon; wt 2; detach=F \\
Econ. & Clip & split \\
Energy & Contract & ren. share $\rightarrow$ CO$_2$/energy; min bundle \\
Energy & Sched. & ep 120; clip=global \\
\bottomrule
\end{tabular*}
\end{table}

\FloatBarrier
\subsection{Data Preprocessing}
\begin{table}[!tbhp]
\centering
\caption{Data preprocessing summary.}
\label{tab:app_data_preprocessing}
\setlength{\tabcolsep}{4pt}
\renewcommand{\arraystretch}{1.05}
\noindent\begin{tabular*}{\columnwidth}{@{\extracolsep{\fill}} l l}
\toprule
Domain & Cleaning \\
\midrule
Synthetic & Balanced; no repair \\
Economics & \(ck, rgdpna>0\); miss.\ $\leq15\%$; lin.\ interp. \\
Energy & ISO3; pop,GDP\(>0\); miss.\ $\leq15\%$; lin.\ interp. \\
\bottomrule
\end{tabular*}
\end{table}


% End
\printbibliography
\vspace{12pt}
\end{document}
